\documentclass[conference]{IEEEtran}

\usepackage[T1]{fontenc}
\usepackage[utf8]{inputenc}
\usepackage{amsmath,amssymb}
\usepackage{graphicx}
\usepackage[numbers,sort&compress]{natbib}
\usepackage{microtype}
\usepackage[hidelinks]{hyperref}
\usepackage{xurl}
\usepackage{float}
\usepackage{tabularx}
\usepackage{array}
\usepackage{placeins}
\usepackage{dblfloatfix}
\newcommand{\resultstableformat}{%
  \footnotesize
  \renewcommand{\arraystretch}{1.12}%
  \setlength{\extrarowheight}{0.7pt}%
}
\setcounter{dbltopnumber}{2}
\renewcommand{\dbltopfraction}{0.95}
\renewcommand{\textfraction}{0.05}


% The report ships with both SVG and PNG figures. Using the PNG copies keeps
% compilation portable and avoids requiring Inkscape or --shell-escape.
\newcommand{\includesvg}[2][]{\includegraphics[#1]{#2.png}}

\title{Where Do Semantics Live? A Homonym-Based Study of Semantic Disambiguation in Large Language Models}

\author{\IEEEauthorblockN{Julia Hagen, 22430997}
\IEEEauthorblockA{Friedrich-Alexander-Universit\"at Erlangen-N\"urnberg}}

\begin{document}
\raggedbottom
\maketitle
\begin{abstract}
Layer-wise comparisons offer one way to test whether artificial and biological
systems organise similar functions across processing stages. We apply this
approach to contextual word meaning in transformers. The study uses seven
two-sense homonyms in four bidirectional encoders and four causal decoders. It
separates three dimensions: layer depth, token position, and incrementally
available context. Sense-conditioned representations are evaluated with
held-out nearest-centroid margins and the Generalized Discrimination Value
(GDV).

Contextual sense was widely decodable, but its strongest readout varied across
depth. Profiling-selected layers reached 0.917 held-out adequacy, compared with
0.893 at the final layer. The crossed model--word interval included zero, and
effects varied across models and homonyms. Adequacy-based cross-word selection
reached 0.914 adequacy, compared with 0.892 for GDV-based selection. The two
criteria also selected different depths in most folds. Global class separation
and item-level decodability therefore captured different properties of the
activation geometry.

Context-position analyses showed that semantic readout also depends on causal
context availability. In causal decoders, homonym-position adequacy was 0.500
when the resolver followed the homonym. At the sentence-final period, after the
resolver had become available, adequacy reached 0.761. Among observations that
were inadequate at the homonym, 67.9\% became locally decodable at the
endpoint.

Incremental context-conflict experiments showed broad updating and persistent
context effects. Resolving information moved a fixed sentinel toward the
authored sense in 77.2\% of observations. Among initially primed states, 40.3\%
crossed into the resolved-sense region. Conflict endpoints remained 0.140 below
matched non-conflicting endpoints containing the same homonym and resolver.

Contextual sense was distributed across depth and depended on the readout
criterion, token position, available context, and context history. These
results motivate functional comparisons between humans and language models at
the level of prediction, integration, conflict, and revision.
\end{abstract}

\begin{IEEEkeywords}
large language models, semantic representations, homonyms, contextual
updating, representational geometry, layer selection, encoders, causal decoders
\end{IEEEkeywords}

\section{Introduction}

Comparisons between artificial and biological information-processing systems
naturally focus on functional analogues: whether both systems perform similar
computations and how those computations are organised internally. Visual
neuroscience provides an established example. Brain-Score evaluates artificial
vision models against neural and behavioural benchmarks and relates model
layers to cortical targets along the primate ventral visual pathway
\citep{schrimpf2020brainscore}. This stage-wise comparison motivates an
analogous question for language: do transformer layers organise linguistic
functions in a reproducible order?

Human comprehension transforms perceptual input into lexical, syntactic, and
semantic representations that support action and language production. These
functions are partly ordered and interact through prediction, bottom-up
evidence, and top-down context \citep{ryskin2023}. Transformers also transform
input through a sequence of internal representations before producing an
output. Their depth may therefore organise linguistic functions across
successive processing stages.

Previous probing studies have found that linguistic information is distributed
non-uniformly across transformer depth \citep{tenney2019}. Different layers can
support different forms of linguistic analysis. The present study tests whether
contextual word sense reaches its most reliable and reproducible readout within
a particular depth region.

This question concerns layer-level population geometry rather than individual
interpretable features. Conceptual features may be distributed or superposed
across many units and layers \citep{elhage2022,bricken2023,templeton2024}, while
their combined expression may become maximally decodable at a narrower stage.

Homonyms provide a controlled assay because the same lexical form can support
sharply different interpretations. If hidden states primarily preserve word
form, occurrences of \emph{bank} should remain similar across financial and
riverside contexts. If contextual sense is expressed, the two groups should
become geometrically distinguishable. The target form is held constant, but
the surrounding sentences still differ in vocabulary, syntax, topic, length,
and target position. The dependent measure therefore captures
\emph{sense-conditioned contextual separability} within the tested contrasts.

\subsection{Locating usable sense information}

A claim that semantics occupies a particular layer depends on how
``semantic information'' is operationalised. This study compares two geometric
criteria. A held-out nearest-centroid margin measures whether an individual
sentence lies closer to the centroid of its authored sense than to the opposing
sense. The Generalized Discrimination Value (GDV) instead summarises labelled
within-sense and between-sense distances across the complete representation
space \citep{schilling2021}. The former emphasises reliable item-level
decisions; the latter emphasises global class geometry without fitting a
classifier.

H1 tests whether selecting a layer from profiling data improves held-out
sense decodability relative to using the final layer. H2 asks whether such a
choice transfers to an unseen homonym and whether margin- and GDV-based
criteria identify the same depth region. Nested and cross-word held-out
evaluation are used because an apparent semantic layer is informative only if
it generalises beyond the observations or word that selected it
\citep{hewitt2019}.

\subsection{Layer depth, token position, and available context}

Layer depth is not equivalent to reading time. Three axes must be distinguished:

\begin{enumerate}
  \item Layer depth indexes successive transformations within one forward pass.
  \item Sequence position determines which context is causally available at a
        particular token.
  \item Incremental input is studied by rerunning the model on progressively
        longer prefixes.
\end{enumerate}

This distinction is especially important when comparing bidirectional encoders
with causal decoders. At a homonym position, an encoder can use both preceding
and following context, whereas a causal decoder can never incorporate a later
resolver there, regardless of depth. A later decoder token can nevertheless
attend to both the homonym and the resolver. Semantic access must therefore be
indexed by both layer and token position.

H3 manipulates whether disambiguating context appears before or after the
homonym and therefore whether it is causally available at that position. H4
then asks whether sense information unavailable at the decoder homonym becomes
locally decodable at a later sentence position after the resolver has been
encountered. Together, these analyses prevent a context-availability effect
from being misinterpreted as a semantic-layer effect.

\subsection{Baseline lean and incremental context conflict}

Even without explicitly disambiguating context, a homonym may lie closer to
one of the two profiled sense regions. H0 measures this baseline asymmetry for
bare words and ambiguous carrier sentences. These states provide a starting
point for interpreting later updating: identical contextual evidence may lead
to different endpoints when representations begin from different positions
along the sense axis.

H5 tests whether that representation changes when new context contradicts an
earlier interpretation. Human comprehension is incremental, and later
disambiguating information can produce processing costs and leave traces of an
initial garden-path interpretation \citep{frazier1982,christianson2001}.
Violation-sensitive neural responses such as the N400 likewise show that
incoming language is evaluated relative to semantic expectations
\citep{ryskin2023,koelbl2026}. These findings motivate a model-internal
context-conflict test. The present
study measures representational change, while matched human data would support
direct behavioural and neural comparison.

Models are therefore rerun on progressively longer prefixes that first prime
one sense and later resolve the homonym toward the other. A fixed sentinel
provides a common readout across stages. The analysis distinguishes movement
toward the resolved-sense region, categorical crossing of the centroid
boundary, and the remaining difference from matched non-conflicting controls.
It provides a model-side measure of geometric updating and persistence.


\subsection{Research questions and analysis structure}
\label{sec:research-questions}

The study addresses three questions.

\begin{description}

\item[\textbf{RQ1}]
\textbf{Where is contextual sense information most accessible?}

Does sense information become especially clear at a reproducible depth, or
does its apparent location vary across models, homonyms, and geometric
criteria?

H1 tests whether an intermediate layer provides a more reliable sense readout
than the final layer. H2 tests whether layer selection transfers across homonyms
and geometric criteria.

\item[\textbf{RQ2}]
\textbf{How does sense accessibility depend on readout position and available
context?}

Does sense decodability depend only on layer depth, or also on the context
available at the selected token position?

H3 compares disambiguating information placed before and after the homonym.
H4 tests whether sense information becomes accessible at a later sentence
position after the resolver has been processed.

\item[\textbf{RQ3}]
\textbf{How flexibly is an established interpretation revised?}

When later context supports a different sense, does the representation move
toward that sense, cross into its geometric region, and approach the endpoint
of a non-conflicting context?

H5 measures incremental updating and the remaining effect of preceding
conflicting context.

\end{description}

H0 supplies the ambiguous-input baseline used to interpret both semantic
localisation and subsequent updating.

Together, these analyses test whether contextual meaning has a reproducible
location in transformer depth and how its representation changes with context.
Agreement across models, homonyms, criteria, and token positions would support
a stable semantic stage. Variation across these dimensions would support a
distributed and path-dependent account.


\section{Data and Methods}
\label{ch:methods}

The study analysed three dimensions of transformer representations. Each
dimension addresses a different source of variation in semantic readout.
Layer depth tests whether sense information becomes most accessible at a
reproducible processing stage. Sequence position tests where the available
context becomes locally readable. Incremental context tests how the readout
changes as new evidence is added across forward passes.

H1--H2 varied layer depth within one forward pass. H3--H4 varied token position
and the context available at that position. H5 varied the supplied context by
rerunning each model on longer prefixes.

\subsection{Shared design}
\label{sec:shared-design}
\label{sec:data}

\paragraph{Operational sense contrasts.}
Seven English homonyms defined seven binary sense contrasts
(Table~\ref{tab:homonyms}). The labels cover the tested senses only.
\emph{Spring} and \emph{pitch}, for example, have additional readings.

Homonyms keep the target form constant across the two sense conditions. This
makes the surrounding context the main systematic source of sense information
at the target. Using seven homonyms tests whether the resulting patterns recur
across lexical items.

An eighth candidate, \emph{light}, was excluded because its intended contrast
paired a noun sense (\emph{illumination}) with an adjective sense
(\emph{low weight}). This would couple sense with grammatical category in the
ambiguity and resolution items.

All stimuli were authored for this study. The homonym remained constant within
each contrast. The surrounding contexts varied in vocabulary, topic, syntax,
length, and sometimes target position. The measures therefore capture
\emph{sense-conditioned contextual separability} and \emph{local sense
decodability} within the seven operational contrasts.

\begin{table}[!t]
\centering
\caption{The seven homonyms and their operational sense contrasts.}
\label{tab:homonyms}
\begin{tabular}{lll}
\hline
Word & Sense 0 & Sense 1 \\
\hline
bank & river / embankment & financial institution \\
bark & dog vocalisation & tree covering \\
bat & sports equipment & flying mammal \\
crane & large wading bird & lifting machine \\
match & sports contest & fire-starting stick \\
spring & season & water source \\
pitch & sports field & business proposal \\
\hline
\end{tabular}
\end{table}

\paragraph{Profiling sentences.}
The profiling sentences defined reference regions for the two senses. Their
unambiguous contexts provided a common geometry against which ambiguous,
reordered, and conflicting inputs could be scored.

The profiling set contained 20 sentences per sense and homonym, giving 280
sentences in total (Table~\ref{tab:h1-bank-stimuli}). Target occurrences were
identified with word-boundary-aware matching. The target could appear in its
singular form or as a regular plural.

When tokenisation split a target into several subwords, their hidden states
were averaged. For example, Qwen2.5 split \emph{Cranes} in ``Cranes migrate
thousands of miles \ldots'' into \texttt{C}, \texttt{ran}, and \texttt{es}.
Their mean represented the target occurrence. These homonym-position states
defined the reference centroids used throughout the study.

For H4, the same profiling pass also retained the final non-special,
non-padding state. This token was the sentence-final period in every profiling
and H4 sentence. Separate period-position centroids allowed H4 to test local
sense decodability at the endpoint in the geometry of that position.

For H5, each profiling sentence was also evaluated with the sentinel
\texttt{probe} appended in a new paragraph. The resulting sentinel states
defined the fixed sense geometry used across the prime, homonym, and resolver
stages.

\begin{table}[!t]
  \centering
  \resultstableformat
  \caption{Sample profiling stimuli for \emph{bank}. The full set contains 20
  sentences per sense. In each H1 outer fold, one sentence is held out. The
  remaining 39 sentences determine the selected layer and the sense centroids.}
  \label{tab:h1-bank-stimuli}

  \begin{tabularx}{\columnwidth}{
    @{}
    >{\centering\arraybackslash}p{0.10\columnwidth}
    >{\raggedright\arraybackslash}p{0.22\columnwidth}
    >{\raggedright\arraybackslash}X
    @{}
  }
    \hline
    Sense & Operational label & Example profiling sentence \\
    \hline
    0 & Riverside &
    There's a steep bank at the bend of the river. \\

    0 & Riverside &
    Fishermen sat quietly on the bank waiting for a bite. \\

    1 & Financial &
    She went to the bank to open a new account. \\

    1 & Financial &
    The bank manager approved her loan application. \\
    \hline
  \end{tabularx}
\end{table}

\paragraph{Models.}
The sample included four bidirectional encoders: ModernBERT-large,
DeBERTa-v3-large, RoBERTa-large, and XLM-RoBERTa-large. It also included four
causal decoders: Qwen2.5-3B, Qwen2.5-7B, Mistral-Nemo-Base-2407, and
OLMo-2-1124-7B. All models were open-weight base checkpoints, which allowed
hidden-state extraction at every layer without instruction tuning.

The sample spans bidirectional and causal attention. This supports the
context-availability question because the two groups differ in access to
following context at the homonym. The broader variation also tests whether
layer and updating patterns recur across model families.

The groups differ in parameter count, training objective, training data,
tokenizer, and depth. Architecture summaries therefore describe the sampled
groups as a whole. The paired context-order manipulation in H3 directly tests
the effect of causal context availability.

\subsection{Representational measures}
\label{sec:representational-measures}

\paragraph{Adequacy margin.}
The research questions require an item-level measure that preserves direction
relative to the authored sense. The adequacy margin provides a common
coordinate for sense classification, margin magnitude, and movement across
positions or context stages.

For a hidden state $h_l$ at layer $l$, let $c_{0,l}$ and $c_{1,l}$ denote the
profiling centroids for the two senses at the relevant readout position. For an
observation with authored sense label $s\in\{0,1\}$, the Euclidean adequacy
margin is

\begin{equation}
M_l(h,s)=
\lVert h_l-c_{1-s,l}\rVert_2
-
\lVert h_l-c_{s,l}\rVert_2.
\label{eq:margin}
\end{equation}

A positive value places the observation closer to its authored-sense centroid.
A negative value places it closer to the opposing centroid. Zero marks the
nearest-centroid boundary. The sign is assigned relative to each observation's
authored sense, so positive values always indicate correct-sense proximity.

The raw margin was normalised by the distance between the two centroids:

\begin{equation}
\widehat{M}_l(h,s)=
\frac{M_l(h,s)}
{\lVert c_{1,l}-c_{0,l}\rVert_2}.
\label{eq:margin-norm}
\end{equation}

The reverse triangle inequality gives
$-1\leq\widehat{M}_l\leq1$. The two centroids attain the endpoint values, and
zero remains the decision boundary. Normalisation reduces the effect of
activation scale across models and layers. Binary adequacy uses the sign
$M_l>0$. Magnitude comparisons use $\widehat{M}_l$.

\paragraph{Centroid estimation and layer use.}
Held-out centroid estimation measures whether the sense geometry transfers to
an unseen observation. H1 and H2 therefore scored each profiling sentence
against centroids estimated from the remaining sentences.

H1 also selected the layer inside each outer leave-one-sentence-out fold.
Layer selection and centroid estimation used only the outer training set. The
outer test sentence therefore measured transfer of both the selected layer and
its reference geometry.

H0 and H3--H5 used centroids and layers estimated from the separate profiling
set. This provided a fixed reference geometry for the ambiguous, reordered,
endpoint, and context-conflict stimuli.

\paragraph{Uncertainty and thresholds.}
The main generalisation claims concern recurrence across both models and
homonyms. Architecture summaries therefore used model--word aggregates as
their evidential units. Contrasts in H1--H3 received 95\% intervals from
20,000 crossed bootstrap resamples of models and homonyms.

H4 reports pooled counts and proportions for its conditional transitions.
H5 reports pooled changes and transition counts, together with consistency
across model--word cells. Its architecture contrasts use 20,000 crossed
bootstrap resamples of models and homonyms. Thresholds such as
$|\widehat{M}_l|>0.3$ appear only in descriptive sensitivity summaries.
Continuous margins and the zero decision boundary remain the primary outcomes.

\subsection{H0: baseline sense lean}
\label{sec:h0}

H0 provides a baseline for RQ1 and RQ3 by locating ambiguous inputs within the
profiled sense geometry. It therefore compared an isolated homonym with five
ambiguous carrier sentences (Table~\ref{tab:h0-bank-stimuli}). The bare input measures
the position produced by the isolated homonym. The carriers test whether
underspecified sentence context produces a stable shift along the same sense
axis.

Each input completed a full forward pass. The homonym state was read at the
independently selected layer for that model--word pair. This keeps the input
condition separate from the depth of the readout.

\begin{table}[!t]
  \centering
  \resultstableformat
  \caption{H0 inputs for \emph{bank}. The bare word and five carriers preserve
  ambiguity between the riverside and financial senses.}
  \label{tab:h0-bank-stimuli}
  \begin{tabular}{p{0.28\columnwidth}p{0.62\columnwidth}}
  \hline
  Condition & Model input \\
  \hline
  Bare word & \texttt{bank} \\
  Ambiguous carrier 1 & The bank was unstable. \\
  Ambiguous carrier 2 & They examined the bank. \\
  Ambiguous carrier 3 & The bank had changed. \\
  Ambiguous carrier 4 & The bank needed protection. \\
  Ambiguous carrier 5 & The bank became difficult to access. \\
  \hline
  \end{tabular}
\end{table}

The two profiling centroids at the same selected layer defined the sense axis.
H1 separately used layer 0 as a diagnostic for surface-form and position
differences.

The signed lean was

\begin{equation}
B_l(h)=
\frac{\lVert h_l-c_{1,l}\rVert_2-\lVert h_l-c_{0,l}\rVert_2}
{\lVert c_{1,l}-c_{0,l}\rVert_2}.
\label{eq:h0-bias}
\end{equation}

Positive values lean toward sense 0. Negative values lean toward sense 1.
This coordinate places the bare and carrier states on the same axis as the two
profiled senses. It therefore shows both the direction and strength of the
baseline asymmetry.

The two label-specific margins for one carrier are exact opposites, so each
carrier contributed one signed state. H0 reports the bare-word lean, the mean
carrier lean, its absolute magnitude, the bare-to-carrier shift, and
carrier-direction consistency. The count with $|B_l|>0.3$ provides a
sensitivity summary.

\subsection{H1: layer depth and held-out adequacy}
\label{sec:h1-method}

A reproducible semantic depth should support unseen sentences from the same
homonym. H1 therefore tested whether a layer selected from profiling sentences
improved held-out sense adequacy relative to the final layer. The final layer
provides a common default that requires no layer selection.

Each profiling sentence was scored at every contextual layer with
leave-one-sentence-out centroids. Selection started at layer 1.

Layer 0 was retained as a surface diagnostic. It precedes self-attention, so
its label separation reflects properties such as token form, tokenisation, and
target position.

Performance used nested leave-one-sentence-out evaluation. Each outer fold
held out one profiling sentence. The remaining sentences selected the layer
with the highest inner leave-one-out adequacy. Mean normalised margin broke
ties. The held-out sentence was then scored at the selected and final layers
with centroids estimated from the outer training set.

A separate full-profile layer was estimated for H0 and H3--H5, which used
independent stimuli. The selected-minus-final contrast measures the practical
value of layer selection. Relative depth was the selected layer index divided
by the index of the final contextual layer. Modal layer and relative depth
measure the stability of the selected location across model--word cells.

\subsection{H2: cross-word layer-selection strategies}
\label{sec:layer-selection}

A general semantic stage should transfer across homonyms. H2 therefore held out
one homonym at a time. The remaining six homonyms selected a layer, and the
held-out homonym measured cross-word transfer with its own
leave-one-sentence-out centroids.

Four strategies were compared:

\begin{enumerate}
    \item \emph{GDV}: the layer with the most negative mean GDV across the six
    profiling homonyms;
    \item \emph{adequacy-based cross-word}: the layer with the highest mean
    leave-one-out adequacy across those homonyms;
    \item \emph{last}: the final layer; and
    \item \emph{oracle}: the best layer for the held-out homonym, used as a
    descriptive ceiling.
\end{enumerate}

The two transferable selectors test whether different geometric objectives
identify the same useful depth. The final layer provides a common default. The
oracle measures the available improvement for the held-out homonym. Ties in
the primary selection score were broken by the mean normalised adequacy
margin, and any remaining tie was resolved by preferring the shallowest tied
layer.

For a layer with $D$-dimensional states, each feature was standardised as
$x'_{nd}=\tfrac{1}{2}(x_{nd}-\mu_d)/\sigma_d$, where $\mu_d$ and $\sigma_d$
were the mean and standard deviation of feature $d$ across that homonym's own
profiling sentences, pooling both senses; a fixed $10^{-12}$ was added to
$\sigma_d$ to avoid division by zero for near-constant features. GDV was then
computed as

\begin{equation}
\operatorname{GDV}=
\frac{\bar d_{\mathrm{intra}}-\bar d_{\mathrm{inter}}}{\sqrt{D}},
\label{eq:gdv}
\end{equation}

where $\bar d_{\mathrm{intra}}$ is the mean within-sense distance and
$\bar d_{\mathrm{inter}}$ is the mean between-sense distance. More negative
values indicate stronger average class separation.

Held-out adequacy measures transfer performance. Oracle regret measures how
much of the available adequacy each strategy recovers. Layer distance and
selection within two layers measure agreement in depth. Crossed model--word
bootstrap intervals compare each transferable strategy with the final layer.

\subsection{H3: causal availability at the homonym position}
\label{sec:h3-method}

RQ2 asks whether homonym-position sense readout depends on access to the
decisive context. H3 therefore placed the same disambiguating clause before or
after the homonym. The two orders preserve the lexical material while changing
whether that clause is available at the homonym in a causal decoder.

Each homonym had five ambiguous carriers. Every carrier appeared with one
clause for each sense and in both clause orders
(Table~\ref{tab:h3-bank-stimuli}). This produced 20 sentences per homonym and
140 sentences in total.

\begin{table}[!t]
    \centering
    \resultstableformat
    \caption{One H3 carrier set for \emph{bank}. Each carrier appears with one
    clause per sense and in both clause orders.}
    \label{tab:h3-bank-stimuli}
    \begin{tabular}{p{0.12\columnwidth}p{0.15\columnwidth}p{0.55\columnwidth}}
    \hline
    Condition & Intended sense & Sentence \\
    \hline
    L & Riverside &
    After weeks of heavy rain, \textbf{the bank was unstable}. \\

    R & Riverside &
    \textbf{The bank was unstable} after weeks of heavy rain. \\

    L & Financial &
    After weeks of market panic, \textbf{the bank was unstable}. \\

    R & Financial &
    \textbf{The bank was unstable} after weeks of market panic. \\
    \hline
    \end{tabular}
\end{table}

The homonym state was read at the profiling-selected layer and scored against
homonym-position centroids. Each L sentence was paired with its R reordering.
This pairing controls for the words in the sentence and isolates clause order.

The primary outcome was the paired L-minus-R change in normalised margin.
Effects were first averaged within model--word cells. The
decoder-minus-encoder interaction tests whether causal access produces a
larger clause-order effect in decoders.

For causal decoders, the two R sentences for one carrier share the same prefix
at the homonym and have opposing labels. Their mean adequacy is therefore 0.500
by construction. The main estimates are the L-minus-R margin and its
architecture interaction.

\begin{table*}[!b]
  \centering
  \resultstableformat
  \caption{One H5 \emph{bank} item. The prime supports the financial sense and
  the resolver selects the riverside sense. Each row is a separate model input.
  The indented \texttt{probe} is the final paragraph of that input.}
  \label{tab:h5-bank-stimuli}

  \begingroup
  \renewcommand{\arraystretch}{1.12}

  \newcommand{\hfiveinput}[1]{%
    \begin{minipage}[t]{\linewidth}
      #1%
      \par\nobreak
      \vspace{0.15em}
      \hspace*{1em}\textbf{\texttt{probe}}
    \end{minipage}%
  }

  \begin{tabular}{
    >{\raggedright\arraybackslash}p{0.17\textwidth}
    >{\raggedright\arraybackslash}p{0.14\textwidth}
    >{\raggedright\arraybackslash}p{0.59\textwidth}
  }
  \hline
  Input condition & Function & Exact input supplied to the model \\
  \hline

  Prime stage &
  Financial prime only &
  \hfiveinput{She finished her loan application and walked to the} \\
  \hline

  Homonym stage &
  Prime plus ambiguous word &
  \hfiveinput{She finished her loan application and walked to the bank} \\
  \hline

  Resolution stage &
  Conflict through the resolver &
  \hfiveinput{She finished her loan application and walked to the bank
  to sit by the river} \\
  \hline

  Matched control &
  Non-conflicting context through the same resolver &
  \hfiveinput{After following the riverside path, she walked to the bank
  to sit by the river} \\
  \hline

  Resolver in isolation &
  Lexical resolver baseline &
  \hfiveinput{river} \\
  \hline
  \end{tabular}

  \endgroup
\end{table*}

\subsection{H4: sequence-position decodability}
\label{sec:h4-method}

RQ2 also asks where sense information becomes readable after later context has
been processed. H4 therefore compared the homonym and sentence-final period
within the same completed forward pass. Both positions were read at the same
selected layer, which keeps depth fixed while token position and available
context change.

The analysis used the ten R-condition sentences per homonym and contained 560
observations.

The homonym state was scored against homonym-position profiling centroids. The
period state was scored against period-position profiling centroids.
Position-specific centroids measure local decodability in the geometry of each
readout. A cross-position diagnostic also scored the period state against the
homonym centroids to test how directly the two geometries transfer.

The gain rate was

\[
P(\text{final adequate}\mid\text{target inadequate}),
\]

and the loss rate was

\[
P(\text{final inadequate}\mid\text{target adequate}).
\]

The gain rate measures how often a target-inadequate observation becomes
adequate under the endpoint-local geometry. The loss rate measures the reverse
change. Because each position has its own centroids, both rates compare local
classification outcomes across positions. H4 reports pooled counts and
proportions for both transitions.

The comparison combines the additional context available at the endpoint with
the change in readout position. Token identity, clause length, and
cue-to-endpoint distance also vary across the two readouts.

\subsection{H5: incremental context-conflict updating}
\label{sec:h5}

RQ3 asks how a sense representation changes when later context supports a new
interpretation. H5 therefore reran each model on progressively longer prefixes.
This produces a separate model state after the prime, homonym, and resolver.

The fixed sentinel provides the same terminal readout token at every stage. The
selected layer and sense axis also remain fixed, which makes the three stages
directly comparable.

The final set contained 98 items, with 14 items per homonym. Seven items per
homonym primed sense 0 and resolved to sense 1. The other seven used the
opposite direction. Each item had a matched non-conflicting control with the
same homonym and resolver. All items followed the order prime, homonym,
resolver.

The sentinel \texttt{probe} was appended in a new paragraph after three
prefixes: the prime, the homonym, and the resolver. The model was rerun for each
prefix. The sentinel state was read at the same profiling-selected layer.
Independent profiling sentences with the same sentinel defined its sense
centroids.

Prefix growth changed the context available to the sentinel and its absolute
sequence position. The stage comparison therefore measures the resulting
change in the terminal readout as context accumulates.

The primary outcome was

\[
\Delta\widehat M=
\widehat M_{\mathrm{resolution}}-
\widehat M_{\mathrm{homonym}}.
\]

A positive value indicates movement toward the authored resolved sense. The
conditional transition rate was

\[
P(\text{resolution correct}\mid\text{homonym primed}).
\]

The resolver-minus-homonym change measures movement toward the resolved sense.
The conditional transition rate measures complete crossing among initially
primed observations. The matched control measures the endpoint cost of
preceding conflict. The resolver-only input measures the sense signal supplied
by the resolver itself. Together, these outcomes separate update magnitude,
boundary crossing, and final endpoint quality.

Thresholds of 0.1 and 0.3 provided sensitivity summaries. The continuous
paired change remained the primary outcome.

A post-hoc architecture comparison first averaged H5 outcomes within each of
the 56 model--word cells. Decoder-minus-encoder differences and 95\% intervals
used 20,000 crossed bootstrap resamples of models and homonyms. These
comparisons describe the sampled architecture groups.

\subsection{Exploratory cross-analysis relationships}
\label{sec:exploratory-methods}

This post-hoc analysis tests whether the H0 starting position predicts later
sense accessibility or updating.

H0 baseline lean was compared with H1 adequacy and H5 updating. The
H0--H1 analysis used all 56 model--word cells. H5 was also separated by
conflict direction, giving 112 model--word--direction cells. Transition rates
were defined for 107 cells. Five cells had an empty primed denominator.
Spearman correlations used these aggregated cells. The analyses were treated as descriptive and report raw correlations.

\FloatBarrier
\clearpage
\section{Results}
\label{ch:results}

\subsection{H0: baseline sense lean}
\label{sec:results-h0}

Across 56 model--word cells, the mean absolute carrier lean was 0.339. Mean
direction consistency across the five carriers was 0.882. The descriptive
threshold $|B_l|>0.3$ was exceeded by 159 of 280 carrier states. Encoders
contributed 65 of 140 cases and decoders contributed 94 of 140
(Table~\ref{tab:h0-arch}).

\begin{table}[h]
\centering
\resultstableformat
\caption{H0 carrier lean by architecture. Mean $|B_l|$ gives the absolute lean
toward either centroid. Direction consistency gives the larger sign proportion
among non-zero carrier leans. The final column counts states with
$|B_l|>0.3$.}
\label{tab:h0-arch}
\begin{tabular}{lccc}
\hline
Architecture & Mean $|B_l|$ & Direction consistency & $|B_l|>0.3$ \\
\hline
Encoder & 0.279 & 0.864 & 65/140 \\
Decoder & 0.398 & 0.900 & 94/140 \\
All models & 0.339 & 0.882 & 159/280 \\
\hline
\end{tabular}
\end{table}

The largest model-specific leans were Qwen2.5-3B toward tree covering for
\emph{bark} ($B_l=-0.657$), Qwen2.5-3B toward water source for
\emph{spring} ($B_l=-0.617$), and OLMo-2-1124-7B toward water source for
\emph{spring} ($B_l=-0.594$). All three cells had complete carrier-direction
agreement.

Several directions also recurred across models. Seven of eight models leaned
toward the financial sense of \emph{bank}. All eight leaned toward tree covering
for \emph{bark}. Seven leaned toward water source for \emph{spring}, and seven
toward the sporting sense of \emph{match}. \emph{Bat} had an across-model mean
near zero and the lowest direction consistency (Figure~\ref{fig:h0-priors}).

Bare-word and carrier directions agreed in 28 of 56 cells. Agreement was
highest for \emph{bank} (7/8 models) and \emph{match} (5/8), and lowest for
\emph{spring} (2/8). Carrier context often moved the mean lean farther from
zero. It could also weaken or reverse the bare-word direction. For
Qwen2.5-7B, bare \emph{crane} strongly favoured bird, while the carriers
consistently favoured machine with a smaller mean magnitude. The contextual
baseline therefore depended on both the homonym and its underspecified carrier.

\begin{figure*}[!t]
\centering
\includesvg[width=\textwidth]{figures/h0_prior_matrix}
\caption{Bare-word and carrier sense lean by model and homonym. Each cell names
the closer sense. Colour saturation shows $|B_l|$. In the carrier panel, a
filled dot marks agreement in at least four of five carriers.}
\label{fig:h0-priors}
\end{figure*}

\begin{table*}[!t]
\centering
\footnotesize
\caption{Carrier audit for the four homonyms with the largest absolute
across-model mean lean. Mean $B_l$ averages eight models and five carriers.}
\label{tab:h0-carrier-audit}
\begin{tabular}{p{0.08\textwidth}p{0.1\textwidth}p{0.09\textwidth}p{0.5\textwidth}}
\hline
Word & Mean direction & Mean $B_l$ & Ambiguous carriers \\
\hline

\emph{bank}
&
Financial
&
$-0.367$
&
The bank was unstable.\newline
They examined the bank.\newline
The bank had changed.\newline
The bank needed protection.\newline
The bank became difficult to access.
\\[0.5em]

\emph{bark}
&
Tree
&
$-0.350$
&
The bark drew attention.\newline
They studied the bark.\newline
The bark was unusual.\newline
The bark changed noticeably.\newline
The bark became the focus.
\\[0.5em]

\emph{spring}
&
Water
&
$-0.329$
&
The spring was unusually dry.\newline
They monitored the spring.\newline
The spring changed quickly.\newline
The spring became a concern.\newline
The spring was hard to assess.
\\[0.5em]

\emph{match}
&
Sport
&
$+0.242$
&
The match was difficult to start.\newline
The match ended quickly.\newline
They checked the match.\newline
The match caused trouble.\newline
The match was watched closely.
\\

\hline
\end{tabular}
\end{table*}

\subsection{H1: nested layer selection}
\label{sec:results-h1}

Layer-0 adequacy was 0.500 across the 56 model--word cells. Word means ranged
from 0.238 for \emph{bank} to 0.675 for \emph{crane}. Because layer 0 precedes
self-attention, these differences reflect surface properties of the profiling
sets. Target position differed between senses for several words. For example,
\emph{bat} followed an average of 4.1 words in sense 0 and 2.0 words in sense
1. The corresponding values for \emph{pitch} were 5.2 and 2.45.

Across 2,240 outer folds, nested selection reached 0.917 adequacy. The final
layer reached 0.893, giving a pooled difference of 2.37 percentage points
(Table~\ref{tab:h1-nested}). The difference was 0.63 points for encoders and
4.11 points for decoders. The crossed model--word effect was 0.0237 with a
95\% interval of $[-0.0058,0.0598]$.

Full-profile layer choices differed by architecture. Encoder choices had a mean
relative depth of 0.762 and a median of 0.866. All 28 encoder--word cells
selected at or above one third of model depth. One cell,
DeBERTa-v3-large on \emph{spring}, selected the one-third boundary. Decoder
choices had a mean relative depth of 0.494 and a median of 0.366. Thirteen of
28 decoder cells selected the first third, and six selected the final layer.
Encoder means were later for every homonym, although the model ranges
overlapped.

Model effects also varied (Figure~\ref{fig:h1-nested}). Selection improved
Qwen2.5-3B by 9.29 points and Mistral-Nemo-Base-2407 by 6.79 points. It reduced
adequacy slightly for Qwen2.5-7B, ModernBERT-large, and RoBERTa-large. The final
layer was therefore a strong default for some models and a weaker choice for
others.

\begin{table}[!t]
\centering
\resultstableformat
\caption{Held-out H1 adequacy under nested layer selection. Difference is
Selected minus Last.}
\label{tab:h1-nested}
\begin{tabular}{lcccc}
\hline
Architecture & Outer folds & Selected & Last & Difference \\
\hline
Encoder & 1,120 & 0.966 & 0.960 & 0.006 \\
Decoder & 1,120 & 0.867 & 0.826 & 0.041 \\
All models & 2,240 & 0.917 & 0.893 & 0.024 \\
\hline
\end{tabular}
\end{table}

\begin{figure*}[!t]
\centering
\includesvg[width=0.77\textwidth]{figures/h1_nested_layer_selection}
\caption{Model-level H1 adequacy at the nested selected layer and final layer.
Labels give the selected-minus-final difference in percentage points.}
\label{fig:h1-nested}
\end{figure*}

\subsection{H2: cross-word selection and the GDV objective}
\label{sec:results-h2}

GDV-selected layers reached 0.892 held-out adequacy, compared with 0.893 at the
final layer (effect $=-0.0009$, 95\% CI $[-0.0366,0.0308]$).
Adequacy-based cross-word selection reached 0.914
(effect $=0.0210$, 95\% CI $[-0.0013,0.0496]$). The held-out-word oracle
reached 0.930 (effect relative to last $=0.0371$, 95\% CI
$[0.0112,0.0728]$) (Table~\ref{tab:h2-strategies}).

Figure~\ref{fig:h2-atlas} gives a separate full-profile comparison within each
model--word cell. It is distinct from the leave-one-word-out transfer analysis.

\begin{figure*}[!t]
\centering
\includesvg[width=0.9\textwidth]{figures/semantic_layer_atlas}
\caption{Full-profile layer choices under nearest-centroid adequacy and GDV.
Each cell shows the shared layer or the two selected layers. The criteria agree
in 26 of 56 model--word cells.}
\label{fig:h2-atlas}
\end{figure*}

\begin{table}[!t]
\centering
\resultstableformat
\setlength{\tabcolsep}{2.5pt}
\renewcommand{\arraystretch}{1.08}
\caption{H2 cross-word layer-selection performance across 56 held-out
model--word folds. Regret and layer distance are measured from the held-out
oracle. Within 2 is the proportion selected within two layers of the oracle.}
\label{tab:h2-strategies}
\begin{tabular}{@{}lcccc@{}}
\hline
Strategy &
Adequacy &
\shortstack{Oracle\\regret} &
\shortstack{Layer\\distance} &
\shortstack{Within\\2 layers} \\
\hline
\noalign{\vskip 0.45ex}
GDV & 0.892 & 0.038 & 4.62 & 0.643 \\
\shortstack[l]{Adequacy-based\\cross-word}
& 0.914 & 0.016 & 6.62 & 0.339 \\
Last layer & 0.893 & 0.037 & 11.96 & 0.250 \\
Oracle & 0.930 & 0.000 & 0.00 & 1.000 \\
\hline
\end{tabular}
\end{table}

The two cross-word selectors chose the same layer in 7 of 56 folds. They chose
layers within two positions in 14 folds, and their median separation was four
layers.

The adequacy difference was concentrated among decoders. Encoder adequacy was
0.960 under GDV, 0.966 under adequacy-based selection, and 0.960 at the final
layer. Their oracle reached 0.980. Decoder adequacy was 0.824 under GDV, 0.862
under adequacy-based selection, and 0.826 at the final layer. Their oracle
reached 0.879.

Transfer also varied by homonym. \emph{Crane} had the lowest adequacy under
GDV-based selection (0.838). \emph{Spring} and \emph{bark} had the largest
oracle regret under GDV, at 0.075 and 0.072. \emph{Bank} was stable under both
criteria (0.900 and 0.903). GDV exceeded adequacy-based selection for
\emph{match} (0.900 versus 0.881) and slightly for \emph{pitch}
(0.888 versus 0.884). For \emph{bat}, the corresponding means were 0.925 and
0.959.

GDV and adequacy summarise different features of a layer. GDV averages
standardised within-class and between-class distances across all observations.
Adequacy tests whether each held-out sentence lies closer to the correct
centroid in the original hidden space. A layer can improve average pairwise
separation while leaving several sentences on the wrong side of the centroid
boundary. Those boundary errors directly reduce adequacy and may have little
effect on GDV. GDV therefore provides a less direct estimate of item-level
decodability.

The Qwen2.5-7B \emph{bat} case shows this difference
(Figures~\ref{fig:h2-geometry} and~\ref{fig:h2-scores}). Layer 26 had the more
negative GDV ($-0.0634$ versus $-0.0505$), while layer 3 had higher
leave-one-out adequacy (0.90 versus 0.60). The PCA projection shows the change
in class geometry. The full-space margins show which sentences crossed the
decision boundary.

\begin{figure*}[!t]
\centering
\includesvg[width=\textwidth]{figures/h2_gdv_geometric_example_overlap_annotated}
\caption{PCA projections of Qwen2.5-7B \emph{bat} states at layers 3 and 26.
The dashed boundary uses the two plotted dimensions. Reported GDV and adequacy
use the full hidden space.}
\label{fig:h2-geometry}
\end{figure*}

\begin{figure*}[!t]
\centering
\includesvg[width=0.9\textwidth]{figures/h2_decision_scores}
\caption{Full-space leave-one-out margins for Qwen2.5-7B on \emph{bat}. Layer
3 assigns 90\% of sentences correctly. Layer 26 assigns 60\% correctly and has
the more favourable GDV.}
\label{fig:h2-scores}
\end{figure*}

\subsection{H3: causal availability at the homonym}
\label{sec:results-h3}

Encoder margins changed little across clause order. Their mean margin was
0.360 in L and 0.358 in R, giving an L-minus-R effect of 0.0026. Decoder
margins changed from 0.342 in L to approximately zero in R, giving an effect
of 0.3420. The decoder-minus-encoder interaction was 0.3394 with a 95\%
interval of $[0.2358,0.4355]$ (Table~\ref{tab:h3-arch}).

\begin{table}[!t]
\centering
\resultstableformat
\caption{H3 margins and adequacy when the decisive clause appears before (L)
or after (R) the homonym.}
\label{tab:h3-arch}
\begin{tabular}{lcccc}
\hline
Architecture & $\widehat M_L$ & $\widehat M_R$ & Adequacy L & Adequacy R \\
\hline
Encoder & 0.360 & 0.358 & 0.843 & 0.861 \\
Decoder & 0.342 & $-0.0002$ & 0.918 & 0.500 \\
\hline
\end{tabular}
\end{table}

The decoder R adequacy of 0.500 confirmed the construction check. All four
decoders showed a large clause-order effect, while each encoder remained near
zero (Figure~\ref{fig:h3-paired}).

\begin{figure*}[!t]
\centering
\includesvg[width=0.77\textwidth]{figures/h3_paired_context_effect}
\caption{Per-model H3 change from L to R. Points show mean
$\widehat M_L-\widehat M_R$ across homonyms. Intervals are 95\%
word-bootstrap intervals.}
\label{fig:h3-paired}
\end{figure*}

The architecture interaction was positive for every homonym. It was largest
for \emph{crane} (0.492), \emph{match} (0.376), and \emph{bat} (0.357), and
smallest for \emph{spring} (0.161). Encoder L-minus-R changes stayed near zero
for most words. \emph{Match} had the largest encoder change at 0.076.

\subsection{H4: local decodability across sequence positions}
\label{sec:results-h4}

Decoder adequacy was 0.500 at the homonym and 0.761 at the sentence-final
period. Among decoder observations that were inadequate at the homonym, 67.9\%
became adequate at the endpoint. Among adequate homonym observations, 15.7\%
became inadequate at the endpoint. Encoder adequacy changed from 0.861 to
0.668, with a gain rate of 41.0\% and a loss rate of 29.0\%
(Table~\ref{tab:h4-transitions}).

\begin{table}[!t]
\centering
\resultstableformat
\setlength{\tabcolsep}{3.5pt}
\renewcommand{\arraystretch}{1.15}
\caption{H4 change in local sense decodability from the homonym to the
sentence-final period. Gain and Loss use the conditional denominators shown.}
\label{tab:h4-transitions}
\begin{tabular}{@{}lcccc@{}}
\hline
Architecture & Target & Final & Gain & Loss \\
\hline
Encoder & 0.861 & 0.668 & 16/39 = 0.410 & 70/241 = 0.290 \\
Decoder & 0.500 & 0.761 & 95/140 = 0.679 & 22/140 = 0.157 \\
\hline
\end{tabular}
\end{table}

Decoder endpoint adequacy exceeded 0.500 for every homonym. Conditional gain
was highest for \emph{crane} (0.950), followed by \emph{pitch} (0.800),
\emph{bank} (0.750), and \emph{bark} (0.750). \emph{Spring} had a gain rate of
0.200 and endpoint adequacy of 0.550. Encoder endpoint adequacy was lower than
homonym adequacy for every word.

The cross-position diagnostic scored period states against homonym centroids.
Its overall adequacy was 0.516. The homonym and period therefore used different
local centroid geometries.

\subsection{H5: fixed-sentinel incremental updating}
\label{sec:results-h5}

The balanced design produced 784 model--item observations. The sentinel mean
was $-0.080$ after the prime, $-0.129$ after the homonym, and $-0.001$ after
the resolver. The resolver-minus-homonym change was $+0.128$. In total,
605/784 observations (77.2\%) moved toward the resolved sense. Among the 580
observations on the primed side at the homonym, 234 (40.3\%) crossed the
boundary by the resolver. Table~\ref{tab:h5-arch} and
Figure~\ref{fig:h5-fixed-sentinel} show the architecture trajectories.

\begin{table*}[!t]
\centering
\footnotesize
\caption{H5 updating by architecture. Prime, Homonym, and Resolver are mean
normalised margins toward the resolved sense. $\Delta$ is Resolver minus
Homonym. Positive $\Delta$ gives the number and proportion that moved toward
the resolved sense. Transition gives the boundary-crossing rate among
observations on the primed side at the homonym.}
\label{tab:h5-arch}
\begin{tabular}{lcccccc}
\hline
Arch. & Prime & Homonym & Resolver & $\Delta$ & Positive $\Delta$ & Transition \\
\hline
\noalign{\vskip 0.51ex}
Encoder & $-0.062$ & $-0.098$ & 0.021 & 0.119 & 284/392 = 72.4\% & 111/262 = 42.4\% \\
Decoder & $-0.098$ & $-0.159$ & $-0.022$ & 0.137 & 321/392 = 81.9\% & 123/318 = 38.7\% \\
\hline
\end{tabular}
\end{table*}

Both architecture groups moved toward the resolved sense. Decoders had the
larger mean change (0.137 versus 0.119), the lower resolver endpoint
($-0.022$ versus 0.021), and the lower transition rate (38.7\% versus 42.4\%).
Their stronger homonym-stage priming placed more observations in the
conditional denominator and farther from the boundary.

The post-hoc decoder-minus-encoder difference in update magnitude was 0.018
with a 95\% interval of $[-0.048,0.080]$. The decoder resolver endpoint was
lower by 0.043 ($[-0.087,-0.001]$). Its conflict-minus-control difference was
lower by 0.084 ($[-0.165,-0.010]$). These estimates describe the sampled
architecture groups.

Decoder updates were larger for five homonyms. Encoders moved farther for
\emph{bark} and \emph{bat}. Encoder resolver endpoints were higher for six
homonyms, while decoders were higher for \emph{pitch}. The mean update was
positive for every model, ranging from 0.074 for ModernBERT to 0.186 for
DeBERTa. Fifty-four of 56 model--word cells had a positive mean update. The
exceptions were XLM-R on \emph{crane} ($-0.009$) and \emph{spring}
($-0.047$).

\begin{figure*}[!t]
\centering
\includesvg[width=0.9\textwidth]{figures/h5_fixed_sentinel_trajectories}
\caption{Fixed-sentinel margin across the prime, homonym, and resolver stages.
Thin lines show model means. Thick lines show encoder and decoder means.}
\label{fig:h5-fixed-sentinel}
\end{figure*}

\begin{table}[!t]
\centering
\resultstableformat
\setlength{\tabcolsep}{2pt}
\setlength{\extrarowheight}{1.5pt}
\renewcommand{\arraystretch}{1.10}

\newcommand{\countpct}[2]{%
  \shortstack[c]{#1\\[0.18em](#2)}%
}

\caption{Word-level H5 outcomes. Mean $\Delta$ is the
resolver-minus-homonym change. Positive $\Delta$ gives the proportion moving
toward the resolved sense. Transition gives the boundary-crossing rate among
initially primed observations. Conflict--control is the resolver-stage margin
minus the matched-control margin.}
\label{tab:h5-word}

\begin{tabular*}{\columnwidth}{
  @{\extracolsep{\fill}}
  lcccc
  @{}
}
\hline
Word &
Mean $\Delta$ &
Positive $\Delta$ &
Transition &
\shortstack{Conflict--\\control} \\
\hline
\noalign{\vskip 0.55ex}

bank   & 0.181 & \countpct{93/112}{83.0\%}
       & \countpct{38/90}{42.2\%} & $-0.162$ \\

bark   & 0.143 & \countpct{101/112}{90.2\%}
       & \countpct{35/83}{42.2\%} & $-0.116$ \\

bat    & 0.122 & \countpct{81/112}{72.3\%}
       & \countpct{31/78}{39.7\%} & $-0.135$ \\

crane  & 0.098 & \countpct{80/112}{71.4\%}
       & \countpct{27/87}{31.0\%} & $-0.174$ \\

spring & 0.068 & \countpct{79/112}{70.5\%}
       & \countpct{24/77}{31.2\%} & $-0.099$ \\

match  & 0.152 & \countpct{97/112}{86.6\%}
       & \countpct{32/82}{39.0\%} & $-0.158$ \\

pitch  & 0.135 & \countpct{74/112}{66.1\%}
       & \countpct{47/83}{56.6\%} & $-0.135$ \\

\hline
\end{tabular*}
\end{table}

Every homonym had a positive mean update. Values ranged from 0.068 for
\emph{spring} to 0.181 for \emph{bank}. Transition rates ranged from about 31\%
for \emph{crane} and \emph{spring} to 56.6\% for \emph{pitch}. Direction also
mattered. Among sense-0-to-sense-1 items, 125/256 initially primed observations
crossed the boundary (48.8\%). The corresponding count for
sense-1-to-sense-0 items was 109/324 (33.6\%).

Several patterns recurred across analyses. \emph{Spring} had the smallest H5
update, the lowest decoder endpoint adequacy in H4, and one of the lowest H5
transition rates. \emph{Crane} had strong H4 recovery and the largest H5
conflict cost. \emph{Bank} combined a strong H0 baseline asymmetry with the
largest H5 update.

The conflict endpoint was 0.140 below the matched control on average. The
difference was $-0.098$ for encoders and $-0.182$ for decoders, and it was
negative in 55 of 56 model--word cells. The conflict endpoint was also 0.056
below the resolver-only baseline. This difference was $-0.018$ for encoders and
$-0.093$ for decoders, and it was negative in 45 of 56 cells
(Figure~\ref{fig:h5-controls}).

\begin{figure*}[!t]
\centering
\includesvg[width=\textwidth]{figures/h5_update_controls}
\caption{Three H5 comparisons at the fixed sentinel. The left panel shows
resolver updating. The middle panel shows conflict cost relative to the matched
control. The right panel compares the conflict endpoint with the resolver-only
baseline.}
\label{fig:h5-controls}
\end{figure*}

Endpoint counts depended on the margin threshold. At zero, 404/784 states
(51.5\%) were on the resolved side. At $|\widehat M|>0.1$, 244 states (31.1\%)
were resolved, 247 (31.5\%) were primed, and 293 (37.4\%) were in the boundary
region. At $|\widehat M|>0.3$, 61 states (7.8\%) were resolved, 57 (7.3\%) were
primed, and 666 (84.9\%) were in the boundary region. The continuous paired
change provides the primary H5 summary.

\subsection{Exploratory relationships across analyses}
\label{sec:results-cross-analysis}

Absolute H0 carrier lean had little association with H1 selected-layer adequacy
across the 56 model--word cells (Spearman $\rho=-0.143$).

For H5, carrier lean was aligned with the primed direction. Stronger primed
alignment predicted a lower correct-sense margin at the homonym
($\rho=-0.227$; 112 cells) and resolver ($\rho=-0.251$; 112 cells). Its
association with resolver-minus-homonym change was near zero
($\rho=-0.026$). Its association with transition rate was weak
($\rho=-0.146$; 107 cells).

This pattern suggests a persistent starting offset. Stronger H0 primed
alignment shifted both the homonym and resolver positions toward the primed
side, while the resolver-induced movement stayed similar. \emph{Bank} combined
a strong baseline lean with the largest update. \emph{Spring} combined a strong
lean with the smallest update. These post-hoc patterns motivate confirmatory
hierarchical analysis.

\FloatBarrier

\section{Discussion}
\label{ch:discussion}

The three research questions produced a consistent account. RQ1 found high
sense decodability but no stable depth across models, homonyms, and geometric
criteria. RQ2 showed that local decodability depends on causal context
availability and readout position. RQ3 showed broad movement toward a resolving
sense alongside partial boundary crossing and a persistent conflict cost.
Together, the results describe contextual sense as distributed,
readout-relative, and path-dependent.

\subsection{Semantic localisation depends on criterion and stimulus}
\label{sec:discussion-localisation}

H0 showed systematic asymmetry in the two-sense geometry before explicit
disambiguation. Some directions matched an intuitive dominant sense, such as
the financial lean of \emph{bank}. Other directions, including tree covering
for \emph{bark} and water source for \emph{spring}, were less obvious.
Bare-word and carrier directions also differed in half of the model--word
cells. The baseline therefore reflects the combined influence of training
distribution, carrier wording, tokenisation, target position, category breadth,
and cluster shape.

These starting positions carried into H5. Stronger alignment with the primed
sense predicted more primed homonym and resolver states. It had little relation
to resolver-induced movement. The models therefore began from unequal
positions while showing similar update magnitudes. This post-hoc pattern
motivates confirmatory testing.

H1 and H2 also showed variable semantic depth. Intermediate layers improved
held-out adequacy for some models, especially the sampled decoders. The mean
advantage varied across model--word cells. GDV-based and adequacy-based
selection often chose different layers. A stable semantic stage would require
greater agreement across models, homonyms, and readout criteria.

The H2 results also clarify the source of the GDV disagreement. GDV summarises
average class geometry. Adequacy tracks held-out decisions at the item level.
A layer can improve average separation while leaving difficult sentences
across the wrong centroid boundary. The Qwen2.5-7B \emph{bat} example shows
this pattern directly. Layer number alone therefore gives an incomplete account
of semantic readout quality.

\subsection{Context availability limits the layer-stage analogy}
\label{sec:discussion-context}

H3 showed that causal availability constrains homonym-position sense readout.
Encoders used right context at the homonym. Causal decoders used only the
preceding prefix. Additional decoder layers transformed that prefix without
access to the later resolver. The resulting clause-order interaction follows
from the attention masks and sets a clear limit on depth-based interpretation.

H4 showed that many decoder cases classified as inadequate at the homonym were
adequate in the endpoint-local geometry. Encoder adequacy followed the opposite
pattern, falling from 0.861 at the homonym to 0.668 at the endpoint despite
full-context access at both positions. Cross-position adequacy was 0.516 when
period states were evaluated against homonym centroids. These results show
strong position specificity: decoder endpoints gained local sense readout,
while encoder endpoints often lost it.

Human comprehension also unfolds incrementally, but it combines recurrence,
prediction, bottom-up input, and top-down context. Causal decoders use a fixed
attention mask, while encoders process full context at once. H3 and H4 therefore
compare access regimes. Their central result is that layer depth, reading time,
and accumulated context form separate axes.

\subsection{Semantic revision is path-dependent}
\label{sec:discussion-updating}

H5 moved from accessibility to revision. The resolver shifted the fixed
sentinel toward the authored sense across all models and homonyms. Movement was
more common than boundary crossing. Many resolver states remained near the
centroid boundary, and 40.3\% of initially primed observations crossed into the
resolved region.

The controls show a lasting effect of context order. Conflict items and matched
controls shared the same homonym and resolver. Conflict endpoints still had
lower correct-sense margins. The resolver-only condition also produced a
strong sense signal. The preceding context shaped how strongly that signal
appeared at the sentinel. Updating and persistence therefore occurred in the
same readout.

Persistence supports continuity while earlier information remains valid.
Sharp contradiction requires stronger revision. In H5, the resolver usually
shifted the representation in the correct direction, while the earlier
interpretation often retained influence. The balance between movement and
boundary crossing indicates contextual inertia under the current geometric
measure.

Human garden-path studies also report traces of an initial interpretation
\citep{christianson2001}. The present study provides a model-side measure of
the same broad function: revision after conflicting context. Matched human
data would allow direct comparison of initial interpretation, revision cost,
and final understanding.

The sampled encoders reached higher resolver endpoints and smaller
conflict--control differences than the sampled decoders. Decoder updates were
sometimes larger because their homonym states began farther toward the primed
sense. Model family, training data, parameter count, tokenisation, and depth
also differed between groups. The architecture summaries therefore describe
the sampled models as a whole.

\subsection{Functional comparison across systems}
\label{sec:discussion-functional}

Brain-Score motivates stage-wise comparison by relating model layers to
biological targets \citep{schrimpf2020brainscore}. For language, the present results motivate a
multidimensional comparison across depth, token position, and context
trajectory.

Behavioural similarity motivates comparisons of computational function. A
human and a transformer may both resolve \emph{bank} from context, while their
intermediate organisation differs. The present results challenge a direct
mapping from transformer depth to a fixed sequence of human-like processing
stages. Semantic readout varied across depth, metric, position, and context
history.

Computational roles offer a clearer basis for comparison. Prediction,
integration, conflict detection, and revision can be measured in both systems.
The N400 is useful because it tracks semantic expectation and violation. The
H5 conflict--control difference provides a model-side conflict measure for a
future matched experiment. Trial-level comparisons could test whether model
signals predict human behaviour or neural activity.

\subsection{Limitations}
\label{sec:discussion-limitations}

The authored sense contexts vary in vocabulary, topic, syntax, target position,
and length. These properties can influence the profiling geometry. Future
stimulus sets should use counterbalanced templates, target masking,
context-only classification, and lexically disjoint train--test splits.

The centroid axis covers two operational senses per homonym. Several words
have additional meanings, and the two profiled regions may differ in breadth
and variability. Adequacy and GDV therefore measure access within the tested
binary contrast. Causal interventions, activation patching, and output measures
would test whether the geometry contributes to model behaviour.

The H5 sentinel fixes token identity and terminal readout role across prefixes.
Longer prefixes also change its absolute position. Length-matched filler
controls, several sentinel tokens, and pooled sentence readouts would separate
these effects and test generality.

The model sample is small and heterogeneous. Repeated observations within
models, homonyms, and items also create dependence beyond the current summary
analyses. Matched model families and hierarchical models would support stronger
architecture claims. Human norms for ambiguity, prime strength, resolver
clarity, reading cost, and final interpretation would support direct
cross-system comparison.

\subsection{Future work}
\label{sec:discussion-future}

A broader stage-order test should apply parallel held-out measures to surface
form, syntax, contextual sense, task instruction, and output-relevant
information. Counterbalanced contrasts could test whether these signals peak
in a stable order, overlap across depth, or reorganise across model families.
Layer studies should also compare broader depth regions.

The relation between H0 and H5 should be tested with a hierarchical model of
baseline lean, prime direction, resolver strength, and update magnitude.
Stronger stimulus controls would separate starting bias from item wording. H5
could also cover corrected entity attributes, changed event states, belief
revision, and repeated updates over longer contexts.

A human experiment should use the same items to measure initial
interpretation, reading or listening cost, final comprehension, and lingering
misinterpretation. Trial-level correspondence would show how closely the model
conflict measure tracks human revision.

\subsection{Conclusion}
\label{sec:discussion-conclusion}

The study found accessible sense structure across models and layers. Its
location varied with the homonym, model, criterion, token position, and context
history. A semantic-stage account must include depth, readout position, and
context history.

Resolving context usually shifted representations toward the authored sense.
Earlier conflicting context still reduced the final correct-sense alignment,
and complete geometric resolution occurred in a minority of initially primed
cases. Semantic readout in these models is distributed, context-sensitive, and
path-dependent.

The results shift the question from locating one semantic layer to explaining
how models make contextual sense information accessible and revise it as
context changes.

\bibliographystyle{IEEEtranN}
\bibliography{references}
\end{document}
